\section{Results}

\subsection*{The support of Stadlmann ($k=49$, claimed $H_1\le240$)}

\begin{theorem}[\lbl{Computed}; exact, two independent programs]\label{thm:240}
For the parameters of \cite{Sta26} in the table above:
\begin{enumerate}
\item for $P$ the cap-free optimiser of degree $20$ in even power sums
($707$ coefficients),
\[49J_T(P1_T)/I_T(P1_T)=4.0120389035\ldots;\]
\item for $P^\ast$ the top eigenvector of $(49G_J,G_I)$ in the same space,
rounded to $400$ significant bits per coefficient,
\[49J_T(P^\ast1_T)/I_T(P^\ast1_T)=4.0124198475\ldots.\]
\end{enumerate}
\end{theorem}

In both cases the two evaluators agree on $I_T$ and $J_T$ as exact rationals.
Without caps the same space gives $4.016176$, so the caps cost $0.0041$ for the
truncated optimiser and $0.0038$ for the optimised one: optimising on $T$
regains less than a tenth of the loss. Since every point of $T$ has all
coordinates at most $0.6$ and $T\subset L\R_{49}$, a dual-weight bound gives
the ceiling $\sup_FkJ_T/I_T\le4.14$ (\lbl{Computed}, interval arithmetic). So
the margin of the variational step of \cite{Sta26} lies in $(0.0124,\,0.14]$.

\begin{corollary}[conditional; the analytic input is \lbl{Claimed}]
The criterion of \cite[Prop.~1]{Sta26} holds for $k=49$ with Stadlmann's
support, caps and parameters. With the \lbl{Lean-checked} admissible
$49$-tuple of diameter $240$, $H_1\le240$ follows from \cite[Prop.~1]{Sta26}
and his equidistribution theorem, which we have not verified.
\end{corollary}

\subsection*{The support of Kintali and of Song--Yue ($k=48$, claimed $H_1\le236$)}
Both preprints use the same support (the table above, called ``D182'' in
\cite{Kin26}) and, by their descriptions, the same public witness of degree
$25$, originally constructed for $k=50$ and evaluated at $k=48$ without
re-optimisation. Song and Yue print it as $1295$ integer triples
$(a,\alpha,c)$, with
\[
P(z)=\sum c\,\left(\tfrac{26}{25}-s(z)\right)^a m_\alpha(z),\qquad
P_U(t)=P\!\left(\tfrac{26}{25U}t\right),\quad U=\tfrac{521}{2000},
\]
$m_\alpha$ the monomial symmetric polynomials, and they print the reference
integers $I_{\mathrm{ref}}=(98!)^2I_0^\ast$ and
$J_{\mathrm{ref}}=(99!)^2(26!)^2J^\ast$. Here $I_0$ is the \emph{cap-free} mass
on $\{s<U\}$, $J$ the capped numerator, and the stars denote integer
coordinates $x=10^4t$. Their criterion is $48J/I_0>2001/2000$, which implies
the capped criterion because $I_T\le I_0$ (\lbl{Literature},
\cite[\S4.12, App.~A, App.~C]{SY26}). They state that a fresh evaluation
reproduces $I_{\mathrm{ref}}$, while fresh regeneration of $J_{\mathrm{ref}}$
from the table ``remains incomplete''.

We extracted the table from the PDF using font sizes to read the exponents in
the partitions. All $1295$ rows are present, every $\alpha$ is nondecreasing
with even parts, and $\max(a+|\alpha|)=25$, as stated. Since all parts are
even, $P_U$ is exactly a polynomial in $L-P_1$ and the even power sums
$P_2,\dots,P_{24}$ ($1311$ terms); the conversion was checked exactly at random
rational points.

\begin{theorem}[\lbl{Computed}, exact rational arithmetic]\label{thm:sy}
For the witness $P_U$ of \cite{SY26} and the support above:
\begin{enumerate}
\item $(98!)^2\,10^{192}\,5^{150}\,521^{50}\,I_0=I_{\mathrm{ref}}$ exactly
(an integer with $655$ digits);
\item $(99!)^2(26!)^2\,10^{196}\,5^{150}\,521^{50}\,J_T=J_{\mathrm{ref}}$
exactly (an integer with $715$ digits), where $J_T$ is the capped numerator;
\item hence $48J_T/I_0=1.000664921591\ldots>2001/2000$, which is
$4.0026596864\ldots$ in the Polymath scaling; with the capped denominator
$48J_T/I_T=4.0030988047\ldots$, and without caps $48J_0/I_0=4.0036348167\ldots$.
\end{enumerate}
\end{theorem}

The factors $10^{192}5^{150}521^{50}$ and $10^{196}5^{150}521^{50}$ come from the
change to integer coordinates when the polynomial is scaled by $65125^{25}$
(our reading of \cite[Eq.~(71)]{SY26}; $65125=25\cdot2605$), so items (1) and
(2) are equalities of integers, not numerical agreement. Item (1) was computed
by a program written for this purpose, which works in the monomial basis of
the table, integrating each product $m_\alpha m_\beta$ by a sum over partial
matchings of parts, and does not use the power-sum conversion. Item (2) was
computed by \texttt{capped\_check.py} on the converted polynomial ($56$ minutes
on four cores). The capped and cap-free values in (3) agree with those claimed
in \cite{Kin26}, $4.0030988$ and $4.0036348$, to every printed digit.

\emph{Status of item (2).} This item and Proposition~\ref{prop:kin-own}
below are the results of this note that so far rest on one program. The second evaluator, \texttt{capped\_exact.py}, agrees
with it on the cap-free part and is computing the $126$ correction pieces
from a checkpoint (about $5$--$10$ hours on four cores). This draft will be
updated when it finishes. The exact agreement of the first program with the
authors' independently produced $715$-digit integer is itself strong evidence
for both, since a common error would have to produce the same integer.

\begin{proposition}[\lbl{Computed}, exact rational arithmetic, one program]\label{prop:kin-own}
For our own cap-free optimiser of degree $22$ for this $(L,\tau)$ at $k=48$
($1041$ coefficients, even power sums, rounded to $400$ significant bits),
$48J_T(P1_T)/I_T(P1_T)=4.0060814208\ldots>4$.
\end{proposition}

Here the caps cost only $0.00033$ (the cap-free value is $4.0064115$). That our
value exceeds the witness's $4.0031$ is expected, since the witness was
optimised for $k=50$ and transported, while ours is optimised for $k=48$.

\begin{corollary}[conditional; the analytic input is \lbl{Claimed}]
The variational criteria of \cite{Kin26} and \cite{SY26} hold at $k=48$, for
the published witness and for an independent test function. With the
\lbl{Lean-checked} admissible $48$-tuple of diameter $236$, $H_1\le236$
follows from their equidistribution input, which we have not verified.
\end{corollary}

\subsection*{The draft of Charton et al.\ ($k=45$, claimed $H_1\le212$)}

\begin{proposition}[\lbl{Computed}; exact, two independent programs]\label{prop:212}
With the caps of \cite[Table~3]{Cha26} and $P$ our cap-free optimiser of
degree $21$ ($846$ coefficients, the dimension of the space
$\mathcal B_{21}$ of \cite[Thm.~11.1]{Cha26}),
$45J_T(P1_T)/I_T(P1_T)=4.0027351128\ldots>4$.
\end{proposition}

The authors report $4.00438$ with their own optimised function (\lbl{Claimed});
our value is lower, as expected for a function optimised without caps, and
already above $4$. The draft, however, also states caps in its Lemma~9.1,
namely $C_1=C_2=31/50$ and $C_m=17/25$ for $m\ge3$, which are smaller than
those of Table~3 for every $m\ge4$. With these the conclusion changes.

\begin{proposition}[\lbl{Computed}, ball arithmetic]\label{prop:ceiling}
Let $T'$ be the support with the caps of Lemma~9.1 of \cite{Cha26} and $V$ the
$846$-dimensional space above. Then $3.99\,G_I-45\,G_J$ is positive definite on
$V$, so
\[
\frac{45J_{T'}(P1_{T'})}{I_{T'}(P1_{T'})}<3.99\qquad\text{for every nonzero }P\in V.
\]
\end{proposition}

The certificate was obtained twice, at $5000$ and $6500$ bits with different
block splits, and the underlying weights agree exactly with
\texttt{capped\_check.py} at a random integer point. The best $P$ in $V$ gives
about $3.98555$ (floating point, from the Gram matrices). So within a space of
the size the authors use, the caps of Lemma~9.1 cannot give their
$4.00438$, which must belong to the caps of Table~3. A careful reading of the
draft supports this: Table~3 is presented as the parameter set used in both
halves of the proof, and the Lemma~9.1 values appear to be left over from an
earlier choice. We report this as a located inconsistency for the authors to
resolve, not as an error in their result. We have not checked that their space
coincides with ours beyond its dimension, and \texttt{capped\_exact.py} has not
yet re-evaluated the random point used for the weights.

